\documentclass[11pt,a4paper]{article}

% ── Geometry ──
\usepackage[margin=1in]{geometry}

% ── Math ──
\usepackage{amsmath,amssymb,amsthm}

% ── Tables ──
\usepackage{booktabs}
\usepackage{array}

% ── Algorithms ──
\usepackage[ruled,vlined,linesnumbered]{algorithm2e}

% ── Hyperlinks ──
\usepackage[colorlinks=true,linkcolor=blue,citecolor=blue,urlcolor=blue]{hyperref}

% ── Misc ──
\usepackage{enumitem}

% ════════════════════════════════════════════════════════════════
% Custom commands
% ════════════════════════════════════════════════════════════════
\newcommand{\Zi}{\mathbb{Z}[i]}
\newcommand{\Gk}{G_k}
\newcommand{\fr}{f(r)}
\newcommand{\Rmoat}{R_{\mathrm{moat}}}
\newcommand{\Rhalf}{R_{0.5}}
\newcommand{\iocount}{\operatorname{io\_count}}
\newcommand{\spant}{\operatorname{span}_T}
\newcommand{\fI}{\mathcal{I}}
\newcommand{\fO}{\mathcal{O}}
\newcommand{\fL}{\mathcal{L}}
\newcommand{\fR}{\mathcal{R}}
\newcommand{\Tplus}{T^+}

% ════════════════════════════════════════════════════════════════
% Theorem environments (numbered by section)
% ════════════════════════════════════════════════════════════════
\theoremstyle{definition}
\newtheorem{definition}{Definition}[section]

\theoremstyle{plain}
\newtheorem{theorem}{Theorem}[section]
\newtheorem{corollary}[theorem]{Corollary}
\newtheorem{lemma}[theorem]{Lemma}

\theoremstyle{remark}
\newtheorem{remark}[theorem]{Remark}

% ════════════════════════════════════════════════════════════════
\title{%
  Gaussian Prime Connectivity Transfer Operator (GPCTO):\\
  Methods for Scalable Moat Detection%
}
\author{---}
\date{%
  2026-03-25\\[0.5em]
  \normalsize\textit{Working methodology document --- Draft v0.2 for peer review}%
}

\begin{document}
\maketitle

% ══════════════════════════════════════════════════════════════
\begin{abstract}
The Gaussian moat problem asks whether, for each fixed step bound~$k$, the connected component of the origin in the Gaussian prime proximity graph~$\Gk$ is finite.
We develop two complementary methods for investigating this question at radii far beyond current computational records.
The \textbf{tile operator} compresses the connectivity of a rectangular region of the Gaussian lattice into boundary data---which connected components touch which faces---and admits a natural composition algebra.
The \textbf{Independent Strip Ensemble (ISE)} deploys many collar-disjoint tile probes at the same radial distance; it is provably conservative (zero false negatives) with false-positive probability decaying exponentially in the number of strips.
For rigorous proof, a \textbf{tile-based upper bound (UB)} method tiles the first-octant annulus gaplessly and checks global disconnection through composition.
We further develop a \textbf{fat-stripe} method---a concrete realization of the UB approach---that tiles an annular strip in the first octant, processes each tile independently via row sieve and sparse union-find, composes all tiles via the GPCTO algebra, and delivers a definitive verdict on whether any connected component spans the annulus radially.
Calibration against the known Tsuchimura moats at $k^2 = 26$ and $k^2 = 32$ confirms correct detection.
At $k^2 = 36$, the ISE locates a percolation transition near $R \approx 80.4\text{M}$; at $k^2 = 40$, results through $R = 3.5\text{B}$ show a sigmoid collapse with $\Rhalf \approx 839\text{M}$.
Fat-stripe calibration confirms total annular barriers at $k^2 = 32$ ($R \approx 2.82\text{M}$) and $k^2 = 36$ ($R \approx 80\text{M}$), with cross-validation against ISE establishing consistency between the two methods.
At $k^2 = 40$, five fat-stripe probes across $R \in [1.03\text{B}, 1.09\text{B}]$ detect total annular blockage---the first such detection at a $k^2$ value where no moat was previously known.  Angular coverage is 128K lattice units (${\sim}\,0.007^\circ$); full-octant tiling is needed for rigorous proof.
\end{abstract}

% ══════════════════════════════════════════════════════════════
\section{Preliminaries}\label{sec:preliminaries}

\subsection{Gaussian integers}

\begin{definition}[Gaussian integers]\label{def:gaussian-integers}
The ring of Gaussian integers is $\Zi = \{a + bi : a, b \in \mathbb{Z}\}$, equipped with the squared norm $|z|^2 = a^2 + b^2$ for $z = a + bi$.
\end{definition}

\subsection{Gaussian primes}

\begin{definition}[Gaussian primes]\label{def:gaussian-primes}
An element $\pi \in \Zi$ is a \emph{Gaussian prime} if it generates a prime ideal in~$\Zi$.  The classification is:
\begin{enumerate}[label=\arabic*.]
  \item \textbf{Split primes.}  If $p \equiv 1 \pmod{4}$ is a rational prime, then $p = a^2 + b^2$ for unique $a > b > 0$ (Cornacchia), and $\pi = a + bi$ and its associates are Gaussian primes with $|\pi|^2 = p$.
  \item \textbf{Inert primes.}  If $p \equiv 3 \pmod{4}$ is a rational prime, then $p$ itself is a Gaussian prime with $|p|^2 = p^2$.
  \item \textbf{Ramified prime.}  The rational prime $2 = -i(1+i)^2$; the element $1+i$ is a Gaussian prime with $|1+i|^2 = 2$.
\end{enumerate}
\end{definition}

We write~$P$ for the set of all Gaussian primes.  By the Gaussian prime number theorem, the count of Gaussian primes with norm at most~$N$ is asymptotically $C \cdot N / \ln N$, where~$C$ depends on convention.  In a lattice rectangle at distance~$R$ from the origin, the local prime density is approximately
\[
  \rho(R) \approx \frac{C}{2\ln R},
\]
with $C \approx 1.274$ calibrated empirically.

\subsection{The Gaussian prime graph}

\begin{definition}[Gaussian prime graph]\label{def:prime-graph}
Fix $k^2 \in \mathbb{N}$.  The \emph{Gaussian prime graph} is $\Gk = (P, E_k)$ where
\[
  \{\pi, \pi'\} \in E_k \quad \Longleftrightarrow \quad |\pi - \pi'|^2 \le k^2.
\]
The \emph{seed set} is $\Sigma_k = \{\pi \in P : |\pi|^2 \le k^2\}$, and the \emph{origin component} $C_0(k)$ is the connected component of~$\Gk$ containing every prime in~$\Sigma_k$.  (One may equivalently adjoin a formal vertex at the origin connected to all primes in~$\Sigma_k$.)
\end{definition}

\subsection{The moat problem}

\begin{definition}[Moat radius]\label{def:moat-radius}
The \emph{moat radius} is
\[
  \Rmoat(k) = \sup\{R \ge 0 : C_0(k) \text{ meets the circle } |z| = R\}.
\]
The Gaussian moat problem asks: is $\Rmoat(k)$ finite for every fixed~$k$?
\end{definition}

\subsection{Prior art}

\begin{table}[ht]
\centering
\begin{tabular}{@{}cll@{}}
\toprule
$k^2$ & $\Rmoat$ & Source \\
\midrule
2  & $\sqrt{137} \approx 11.7$           & Classical \\
6  & $\approx 3{,}612$                   & Tsuchimura 2005 \\
10 & $\approx 26{,}862$                  & Tsuchimura 2005 \\
18 & $\approx 175{,}045$                 & Tsuchimura 2005 \\
26 & $\approx 1{,}015{,}639$             & Tsuchimura 2005 \\
32 & $\approx 2{,}823{,}055$             & Tsuchimura 2005 \\
36 & Unknown (ISE transition at $\approx 80\text{M}$)  & This work \\
40 & Unknown (ISE transition at $\approx 839\text{M}$) & This work \\
\bottomrule
\end{tabular}
\caption{Known and estimated moat radii.}
\label{tab:prior-art}
\end{table}

Tsuchimura's method grows the origin component outward through all Gaussian primes, at cost $O(R^2 / \ln R)$ in both time and memory.  This becomes prohibitive for $R > 10^8$.


% ══════════════════════════════════════════════════════════════
\section{The Tile Operator}\label{sec:tile-operator}

\subsection{Tile definition}

\begin{definition}[Tile]\label{def:tile}
A \emph{tile} is a lattice rectangle
\[
  T = T(a_0, b_0, W, H) = \{a + bi \in \Zi : a_0 \le a \le a_0 + H,\ b_0 \le b \le b_0 + W\},
\]
where~$a$ is the radial coordinate and~$b$ the lateral coordinate.  We call~$W$ the \emph{width} and~$H$ the \emph{height} of the tile.
\end{definition}

\subsection{The four faces}

\begin{definition}[Faces]\label{def:faces}
The tile~$T$ has four distinguished faces.  A Gaussian prime $\pi = a + bi \in P \cap T$ touches:
\begin{itemize}
  \item the \emph{Inner face}~$\fI$ if $a - a_0 \le c$,
  \item the \emph{Outer face}~$\fO$ if $(a_0 + H) - a \le c$,
  \item the \emph{Left face}~$\fL$ if $b - b_0 \le c$,
  \item the \emph{Right face}~$\fR$ if $(b_0 + W) - b \le c$,
\end{itemize}
where~$c$ is the collar width defined below.  These inequalities are non-strict ($\le$, not $<$): a prime at distance exactly~$c$ from a face boundary is counted as touching that face.
\end{definition}

\subsection{Collar}

\begin{definition}[Collar]\label{def:collar}
The \emph{collar} is
\[
  c = \left\lceil \sqrt{k^2} \right\rceil.
\]
The \emph{expanded tile} is
\[
  \Tplus = \{a + bi : a_0 - c \le a \le a_0 + H + c,\ b_0 - c \le b \le b_0 + W + c\}.
\]
\end{definition}

\noindent\textbf{Justification.}  An integer step vector $(\Delta a, \Delta b)$ satisfying $\Delta a^2 + \Delta b^2 \le k^2$ has maximal single-coordinate excursion
\[
  \max\{|\Delta a| : \exists\,\Delta b,\ \Delta a^2 + \Delta b^2 \le k^2\} = \left\lfloor \sqrt{k^2} \right\rfloor.
\]
The floor is the tight bound: the maximum~$|\Delta a|$ is achieved at $\Delta b = 0$, giving $|\Delta a| \le \lfloor\sqrt{k^2}\rfloor$.

Using $c = \lceil\sqrt{k^2}\rceil$ rather than $\lfloor\sqrt{k^2}\rfloor$ is the \emph{conservative} choice.  When~$k^2$ is a perfect square, $\lceil\sqrt{k^2}\rceil = \lfloor\sqrt{k^2}\rfloor$, so the two definitions coincide.  When~$k^2$ is not a perfect square, $\lceil\sqrt{k^2}\rceil = \lfloor\sqrt{k^2}\rfloor + 1$, so the collar is one unit wider than strictly necessary.

This over-collaring is harmless for correctness---the expanded tile includes strictly more lattice points than required---and it provides two benefits:
\begin{enumerate}
  \item \textbf{Uniform formula.}  For any~$k^2$, the collar $c = \lceil\sqrt{k^2}\rceil$ guarantees that every edge incident to an interior prime is visible in~$\Tplus$, without case analysis on whether~$k^2$ is a perfect square.
  \item \textbf{Face-port completeness.}  With~$\le$ in the face-membership inequalities, a prime at distance exactly~$c$ from the boundary is tagged as a face port.  Using ceil ensures that this threshold is never below the actual maximum excursion, even though the excess is at most one unit.
\end{enumerate}

\noindent\emph{Concrete values:}

\begin{table}[ht]
\centering
\begin{tabular}{@{}cccc@{}}
\toprule
$k^2$ & $\lfloor\sqrt{k^2}\rfloor$ & $\lceil\sqrt{k^2}\rceil$ & Remark \\
\midrule
36 & 6 & 6 & Perfect square; floor = ceil \\
40 & 6 & 7 & Non-perfect square; collar oversized by 1 \\
\bottomrule
\end{tabular}
\caption{Collar values for key step bounds.}
\label{tab:collar-values}
\end{table}

Throughout this document, $c = \lceil\sqrt{k^2}\rceil$ is the canonical collar.  All theorems are stated and proved using this definition.

\subsection{Face ports}

\begin{definition}[Face ports]\label{def:face-ports}
A \emph{face port} is a record $(\pi, C, F)$ where $\pi \in P \cap T$ is a Gaussian prime in the interior tile, $C$~is the index of its connected component in~$\Gk[\Tplus]$, and $F \in \{\fI, \fO, \fL, \fR\}$ is the face it touches.  A single prime may be a face port for multiple faces.
\end{definition}

\subsection{Tile connectivity function}

\begin{definition}[Tile connectivity function]\label{def:tile-connectivity}
Given tile~$T$ and step bound~$k^2$:
\begin{enumerate}
  \item Form the induced subgraph $\Gk[\Tplus]$ on all Gaussian primes in the expanded tile~$\Tplus$.
  \item Compute the connected components of $\Gk[\Tplus]$ via union-find.
  \item For each component~$C$, determine its \emph{face set} $\partial_T C \subseteq \{\fI, \fO, \fL, \fR\}$---the set of faces touched by primes of~$C$ that lie in the interior tile~$T$.
  \item The \emph{transfer data} are the pairwise face-spanning counts:
  \[
    \spant(X, Y) = \#\{C : X, Y \in \partial_T C\}, \qquad X, Y \in \{\fI, \fO, \fL, \fR\},\ X \ne Y.
  \]
\end{enumerate}
\end{definition}

\subsection{Inner-to-Outer count}

\begin{definition}[\texorpdfstring{$\iocount$}{io\_count}]\label{def:io-count}
The \emph{inner-to-outer count} is
\[
  \iocount(T) = \spant(\fI, \fO),
\]
the number of distinct connected components in~$\Gk[\Tplus]$ that touch both the Inner and Outer faces of~$T$.
\end{definition}

\noindent\textbf{Interpretation.}  $\iocount(T) > 0$ means there exists a path in~$\Gk[\Tplus]$ from a prime near the bottom edge to a prime near the top edge.  $\iocount(T) = 0$ means the tile is \emph{radially blocked}: no connectivity passes through it in the outward direction.


% ══════════════════════════════════════════════════════════════
\section{Tile Composition}\label{sec:composition}

\subsection{Horizontal composition}

\begin{definition}[Horizontal composition]\label{def:h-composition}
Let $T_L = T(a_0, b_0, W, H)$ and $T_R = T(a_0, b_0 + s_b, W, H)$ be two tiles at the same radial level, separated laterally by stride~$s_b$.  Their \emph{horizontal composition} $T_L \otimes_h T_R$ is a new tile operator whose components are obtained by:
\begin{enumerate}
  \item Taking the union of all components from both tiles.
  \item Merging component~$C_L$ from~$T_L$ with component~$C_R$ from~$T_R$ whenever there exist face ports $p_L \in \fR(T_L)$ and $p_R \in \fL(T_R)$ belonging to~$C_L$ and~$C_R$ respectively, with $|p_L - p_R|^2 \le k^2$.
  \item Recomputing face sets for the merged components.
\end{enumerate}
The composed tile has faces $\fI(T_L) \cup \fI(T_R)$, $\fO(T_L) \cup \fO(T_R)$, $\fL(T_L)$, and $\fR(T_R)$.
\end{definition}

\subsection{Vertical composition}

\begin{definition}[Vertical composition]\label{def:v-composition}
Let $T_B = T(a_0, b_0, W, H)$ and $T_T = T(a_0 + H, b_0, W, H)$ be two tiles stacked radially.  Their \emph{vertical composition} $T_B \otimes_v T_T$ merges components whenever there exist face ports $p_B \in \fO(T_B)$ and $p_T \in \fI(T_T)$ with $|p_B - p_T|^2 \le k^2$.  The composed tile has faces $\fI(T_B)$, $\fO(T_T)$, $\fL(T_B) \cup \fL(T_T)$, and $\fR(T_B) \cup \fR(T_T)$.
\end{definition}

\subsection{Composition soundness}

\begin{theorem}[Distance-based composition is conservative]\label{thm:composition-sound}
Let~$T_1$, $T_2$ be tiles sharing a boundary, and let $T_1 \otimes T_2$ denote their composition via distance-based face-port matching.  If two components~$C_1$, $C_2$ are merged in $T_1 \otimes T_2$, then there exists a genuine edge $\{\pi_1, \pi_2\} \in E_k$ with $\pi_1 \in C_1$, $\pi_2 \in C_2$.
\end{theorem}

\begin{proof}
The merge occurs because there exist face ports $p_1 \in C_1$ and $p_2 \in C_2$ with $|p_1 - p_2|^2 \le k^2$.  By Definition~\ref{def:prime-graph}, this means $\{p_1, p_2\} \in E_k$.  Since $p_1 \in C_1$ and $p_2 \in C_2$, the edge connects the two components in~$\Gk$.
\end{proof}

\begin{corollary}\label{cor:no-false-connections}
Distance-based composition never introduces false connections: every merged pair of components is genuinely connected in~$\Gk$.
\end{corollary}

\begin{remark}[Conservative direction]\label{rem:conservative}
Distance-based composition can, however, \emph{miss} connections.  Two primes $\pi_1 \in T_1$ and $\pi_2 \in T_2$ might be connected in~$\Gk$ via a multi-hop path that passes through primes in the overlap region without either~$\pi_1$ or~$\pi_2$ being face ports.  An alternative \emph{shared-prime matching} scheme---identifying primes with identical coordinates in the overlap zone and merging their respective component assignments---captures all such connections exactly.  The current implementation uses distance-based matching; shared-prime matching is formalized but not yet implemented.  For the ISE method (Section~\ref{sec:ise}), composition is not used and this distinction is immaterial.  For the tile-UB method (Section~\ref{sec:ub}), distance-based matching is sound (it can only undercount connections, which is conservative for proving disconnection) but potentially less tight than shared-prime matching.
\end{remark}


% ══════════════════════════════════════════════════════════════
\section{Independent Strip Ensemble (ISE)}\label{sec:ise}

\subsection{Strip definition}

\begin{definition}[Strip]\label{def:strip}
A \emph{strip}~$S_j$ at radial level~$R$ with tile dimensions $W \times H$ is a sequence of tiles:
\[
  S_j = \bigcup_{n=0}^{N-1} T(a_{\text{lo}} + nH,\; b_j,\; W,\; H),
\]
where $a_{\text{lo}}$ is chosen to place the inner face at radial distance~$R$ from the origin, $b_j$~is the lateral offset of strip~$j$, and $N = \lceil(R_{\max} - R_{\min})/H\rceil$ is the number of shells.
\end{definition}

\begin{definition}[Independence criterion]\label{def:independence}
Two strips~$S_j$ and~$S_l$ are \emph{independent} if their expanded tiles are disjoint in the lateral coordinate: the collar-expanded regions $[b_j - c,\; b_j + W + c]$ and $[b_l - c,\; b_l + W + c]$ do not overlap.  This is guaranteed when the lateral stride satisfies
\[
  |b_j - b_l| \ge W + 2c.
\]
The standard ISE layout places~$M$ strips at positive-only offsets:
\[
  b_j = c + j \cdot (W + 2c), \qquad j = 0, 1, \ldots, M-1.
\]
This ensures all~$M$ expanded tiles are laterally disjoint and avoids the $b \to -b$ reflection symmetry (Gaussian primality of $a + bi$ is identical to that of $a - bi$), which would halve the effective independent sample count.
\end{definition}

\subsection{The \texorpdfstring{$\fr$}{f(r)} metric}

\begin{definition}[$\fr$]\label{def:fr}
At radial shell~$n$ with midpoint $r_n = R_{\min} + (n + \tfrac{1}{2})H$, define
\[
  f(r_n) = \frac{\#\{j : \iocount_j(r_n) > 0\}}{M},
\]
where $\iocount_j(r_n)$ is the inner-to-outer count of the tile in strip~$j$ at shell~$n$.
\end{definition}

\noindent\textbf{Interpretation.}  $\fr$ is the fraction of independent strip probes that detect Inner-to-Outer connectivity at radius~$r$.  It ranges from~$0$ (all strips blocked) to~$1$ (all strips crossing).

\subsection{Subgraph monotonicity (zero false negatives)}

\begin{theorem}[Subgraph monotonicity]\label{thm:monotonicity}
Let $S \subseteq \Zi$ be any subset containing the seed set~$\Sigma_k$.  Let $R_{\mathrm{strip}}(k, S)$ denote the moat radius of the origin component in the induced subgraph~$\Gk|_S$.  Then
\[
  R_{\mathrm{strip}}(k, S) \le \Rmoat(k).
\]
\end{theorem}

\begin{proof}
The primes in~$S$ form a subset $P_S \subseteq P$.  The graph $\Gk|_S = (P_S, E_k|_{P_S})$ is an induced subgraph of~$\Gk$.  Every path in~$\Gk|_S$ is also a path in~$\Gk$ (each edge satisfies $|\pi - \pi'|^2 \le k^2$ and both endpoints lie in~$P$).  Therefore every prime reachable from~$\Sigma_k$ in~$\Gk|_S$ is also reachable in~$\Gk$.  Taking suprema of radii gives $R_{\mathrm{strip}}(k, S) \le \Rmoat(k)$.
\end{proof}

\begin{corollary}[Local crossing soundness]\label{cor:local-crossing}
If a tile~$T$ has $\iocount(T) > 0$, then there exists a genuine Inner-to-Outer path in~$\Gk$ through the collared region~$\Tplus$.
\end{corollary}

\begin{proof}
The component witnessing $\iocount(T) > 0$ lives in~$\Gk[\Tplus]$, which is a subgraph of~$\Gk$.  The same path exists in~$\Gk$.
\end{proof}

\begin{theorem}[ISE zero false negatives]\label{thm:ise-zero-fn}
Let $\{S_j\}_{j=0}^{M-1}$ be an ensemble of strips covering a radial interval $[R, R+H]$.  If $\fr = 0$ at every strip, then no path in~$\Gk$ passes from $|z| < R$ to $|z| > R + H$ through any of the~$M$ sampled lateral positions.
\end{theorem}

\begin{proof}
For each strip~$S_j$, the tile at shell~$r$ computes $\Gk[\Tplus_j]$, which contains every Gaussian prime in the expanded tile and every edge incident to any interior prime.  If $\iocount_j(r) = 0$, then no connected component of~$\Gk[\Tplus_j]$ spans from the Inner to the Outer face of~$T_j$.

Any path in~$\Gk$ from a prime with $a < a_{\text{lo}} + c$ to a prime with $a > a_{\text{lo}} + H - c$, passing through the lateral range $[b_j, b_j + W]$, must consist entirely of primes within~$\Tplus_j$ (since each step moves at most~$c$ units laterally, and the strip width~$W$ exceeds~$2c$, a path entering the strip's interior cannot exit the expanded tile's lateral range and re-enter within the same tile height).  Therefore any such path would appear as a component in~$\Gk[\Tplus_j]$, contradicting $\iocount_j = 0$.
\end{proof}

\begin{remark}[False positives are possible]\label{rem:false-positives}
$\iocount(T) > 0$ certifies a local crossing but does not imply the crossing is part of the origin component~$C_0(k)$.  The ISE measures \emph{local percolation} (whether any component crosses a shell), not \emph{origin-connected reachability} (whether $C_0(k)$ crosses a shell).  These are distinct: a Tsuchimura moat can exist at a radius where $\fr > 0$, if the crossing components are disconnected from the origin.
\end{remark}

\subsection{False-positive analysis}

\begin{theorem}[Ensemble false-positive bound]\label{thm:fp-bound}
Suppose, at some shell~$r$, the full graph~$\Gk$ contains an Inner-to-Outer crossing.  Let $p(r)$ denote the probability that a single strip fails to detect this crossing (i.e., $\iocount_j(r) = 0$ despite the crossing existing in~$\Gk$).  If the~$M$ strips are independent (collar-disjoint), then
\[
  \Pr(\fr = 0 \mid \text{crossing exists in } \Gk) \le p(r)^M.
\]
\end{theorem}

\begin{proof}
The event $\fr = 0$ requires all~$M$ strips to independently fail.  Under the disjointness condition, the prime populations in distinct tiles are disjoint, so their $\iocount$ outcomes are determined by disjoint subsets of primes.  The product bound follows.
\end{proof}

Over~$N$ shells, the union bound gives $\Pr(\text{any false positive}) \le N \cdot p(r)^M$.  Concrete values:

\begin{table}[ht]
\centering
\begin{tabular}{@{}cccc@{}}
\toprule
$p$ & $M = 16$ & $M = 32$ & $M = 64$ \\
\midrule
0.30 & $4.3 \times 10^{-9}$  & $1.9 \times 10^{-17}$ & $3.5 \times 10^{-34}$ \\
0.10 & $10^{-16}$            & $10^{-32}$            & $10^{-64}$ \\
0.01 & $10^{-32}$            & $10^{-64}$            & $10^{-128}$ \\
\bottomrule
\end{tabular}
\caption{False-positive probability $p^M$ for various single-strip failure probabilities.}
\label{tab:fp-bounds}
\end{table}

\begin{remark}[Independence model]\label{rem:independence}
The product formula requires that $\iocount$ outcomes across strips are independent.  This holds exactly when the expanded tiles are laterally disjoint (guaranteed by stride $\ge W + 2c$).  Gaussian primes are deterministic, so ``independence'' here means that the prime distribution in one tile is not predictive of the connectivity pattern in a distant tile.  For tiles separated by $\gg \sqrt{k^2}$, this is empirically well-supported: the Hardy--Littlewood conjecture for Gaussian primes predicts asymptotic independence of prime counts in disjoint regions.
\end{remark}

\subsection{Curvature compensation}

\noindent\textbf{Problem.}  Strips are rectangular, but the moat is circular.  Under naive placement (all strips sharing the same $a_{\text{lo}}$), a strip at lateral offset~$b_j$ has its Inner face center at Euclidean distance $\sqrt{a_{\text{lo}}^2 + (b_j + W/2)^2}$ from the origin.  Lateral strips are therefore at slightly different radii.

\begin{definition}[Compensated placement]\label{def:compensated}
For strip~$j$ with lateral center $b_{\text{center},j} = b_j + W/2$, set
\[
  a_{\text{lo},j} = \left\lfloor \sqrt{R_{\text{target}}^2 - b_{\text{center},j}^2} \right\rfloor,
\]
where $R_{\text{target}}$ is the desired inner-face radius for all strips.
\end{definition}

\begin{theorem}[Compensated placement accuracy]\label{thm:compensated}
Under compensated placement, the actual inner-face center distance from the origin satisfies
\[
  \left|\sqrt{a_{\text{lo},j}^2 + b_{\text{center},j}^2} - R_{\text{target}}\right| < 1.
\]
\end{theorem}

\begin{proof}
By definition, $a_{\text{lo},j} = \lfloor x \rfloor$ where $x = \sqrt{R_{\text{target}}^2 - b_{\text{center},j}^2}$.  Then $x - 1 < a_{\text{lo},j} \le x$, so
\[
  a_{\text{lo},j}^2 + b_{\text{center},j}^2 \le x^2 + b_{\text{center},j}^2 = R_{\text{target}}^2,
\]
and
\[
  a_{\text{lo},j}^2 + b_{\text{center},j}^2 > (x-1)^2 + b_{\text{center},j}^2 = R_{\text{target}}^2 - 2x + 1.
\]
Since $x > 0$, the squared distance lies in $(R_{\text{target}}^2 - 2x + 1,\; R_{\text{target}}^2]$, giving a radial error bounded by $1/R_{\text{target}} < 1$ for $R_{\text{target}} > 1$.
\end{proof}

\begin{remark}\label{rem:compensated-independence}
Compensated placement preserves strip independence: only the $a$-coordinate changes per stripe; the lateral separation $|b_j - b_l| \ge W + 2c$ is unchanged.
\end{remark}

\subsection{Angular coverage and stripe count}

The ISE covers an angular sector of the first quadrant from $\theta = 0$ (the positive real axis) to $\theta_{\max} = \arctan(b_{\max} / R)$, where $b_{\max} = c + (M-1)(W + 2c) + W$ is the rightmost lateral extent.

For $M = 32$, $W = 2000$, $c = 7$ (at $k^2 = 40$): $b_{\max} = 7 + 31 \cdot 2014 + 2000 = 64{,}441$.  At $R = 1\text{B}$: $\theta_{\max} \approx 0.0037^\circ$.  The angular coverage is a tiny fraction of the full annulus, which is precisely the point: the ISE does not need angular coverage.  It needs~$M$ independent samples, and by the $p^M$~bound, even $M = 32$ collar-disjoint samples suffice for exponentially reliable detection.


% ══════════════════════════════════════════════════════════════
\section{Tile-Based Upper Bound (UB) Method}\label{sec:ub}

\subsection{Concept}

The ISE detects moat \emph{candidates} but cannot rigorously prove a moat exists: it samples only a sparse set of angular positions and measures local percolation rather than origin-connected reachability.  To \emph{prove} $\Rmoat(k) \le R$ for some~$R$, one must demonstrate that no path in~$\Gk$ crosses from $|z| < R$ to $|z| > R + H$ through any angular position.

\subsection{Octant symmetry}

The dihedral symmetry group~$D_4$ of~$\Zi$---generated by the quarter-turn $z \mapsto iz$ and conjugation $z \mapsto \bar{z}$---preserves Gaussian primality.  Consequently, it suffices to verify disconnection in the \emph{first octant}: the region $\{a + bi : 0 \le b \le a\}$.

\begin{lemma}[Octant barrier implies full-plane moat]\label{lem:octant}
If no connected component of~$\Gk$ restricted to the first-octant rectangular band $\{a + bi : R_{\min} \le a \le R_{\max},\; 0 \le b \le a\}$ has a component crossing from $a = R_{\min}$ to $a = R_{\max}$, then $\Rmoat(k) \le R_{\max}\sqrt{2}$.
\end{lemma}

\begin{proof}[Proof sketch]
Any path from the origin to $|z| > R_{\max}\sqrt{2}$ must eventually have $\max(|a|, |b|) > R_{\max}$, hence $a > R_{\max}$ in at least one of the eight octants (by the octant symmetries).  The rectangular barrier in the first octant, transported to all eight octants by~$D_4$, blocks any such path.
\end{proof}

\begin{remark}\label{rem:octant-conservative}
The first-octant restriction is conservative for the UB\@: connections crossing the octant boundaries ($b = 0$ and $b = a$) are ignored, which can only undercount connectivity.
\end{remark}

\subsection{Gapless tiling geometry}

\begin{definition}[Gapless tiling]\label{def:gapless}
A \emph{gapless tiling} of the first-octant annulus at radius~$R$ is a collection of tiles $\{T_{n,j}\}$ such that:
\begin{enumerate}
  \item Every Gaussian prime~$\pi$ with $R_{\min} \le \mathrm{Re}(\pi) \le R_{\max}$ and $0 \le \mathrm{Im}(\pi) \le \mathrm{Re}(\pi)$ lies in at least one tile's interior.
  \item Tiles in the same radial shell~$n$ are spaced with lateral stride $s_b = W$ (not $W + 2c$), so tile interiors partition the band laterally.
  \item Adjacent tiles share expanded regions of width~$2c$ in the overlap zone.
\end{enumerate}
\end{definition}

\begin{definition}[Tiles per radial shell]\label{def:tiles-per-shell}
At radius~$R$, the first-octant lateral extent is approximately~$R$ (from $b = 0$ to $b = a \approx R$).  The number of tiles per shell is
\[
  N_\ell = \left\lceil \frac{R}{W} \right\rceil.
\]
At $R = 1.5\text{B}$, $W = 2000$: $N_\ell = 750{,}000$.
\end{definition}

\subsection{Composition for UB}

Tiles in each radial shell are composed horizontally (Left/Right matching) to produce a \emph{band operator} for the entire first-octant width.  Then consecutive radial shells are composed vertically (Inner/Outer matching).  If the final composed operator has $\iocount = 0$ for the full octant width, the moat is confirmed at radius~$R$.

The composition uses distance-based face-port matching (Theorem~\ref{thm:composition-sound}), which is sound for the UB direction: it can only undercount connections, making ``blocked'' verdicts reliable.

\subsection{ISE-to-UB handoff}

The ISE provides two quantities that guide the UB campaign:
\begin{enumerate}
  \item $\Rhalf$: the radius where $\fr = 0.5$ (the percolation transition midpoint).
  \item $R_{\text{ext}}$: the radius where $\fr \approx 0$ (the extinction radius).
\end{enumerate}
The UB campaign starts at~$R_{\text{ext}}$, where ISE predicts disconnection is most likely.  This avoids a blind search from the origin: instead of scanning all radii, start where the signal indicates the moat is.

\subsection{Computational feasibility}\label{sec:ub-feasibility}

\subsubsection{Current \texttt{tile\_cuda} architecture}

The \texttt{tile\_cuda} binary splits tile processing across GPU and CPU\@.  Measured on a Jetson Orin Nano with $2000 \times 2000$ tiles at $k^2 = 36$:

\begin{table}[ht]
\centering
\begin{tabular}{@{}llr@{}}
\toprule
Stage & Device & Time \\
\midrule
GPU primality (sieve + Miller--Rabin)  & GPU & ${\sim}\,170\text{ ms}$ \\
Device-to-host bitmap copy             & DMA & ${\sim}\,2\text{ ms}$ \\
CPU union-find + face classification   & CPU & ${\sim}\,40\text{ ms}$ \\
CUDA context init (one-time)           & GPU & ${\sim}\,150\text{ ms}$ \\
\midrule
\textbf{Total per tile} (after init)   &     & $\mathbf{{\sim}\,330\textbf{ ms}}$ \\
\bottomrule
\end{tabular}
\caption{Per-tile timing breakdown on Jetson Orin Nano (\texttt{tile\_cuda}).}
\label{tab:tile-cuda-timing}
\end{table}

\subsubsection{GPU union-find gap}

The current implementation runs union-find on CPU\@.  The \texttt{gpu\_cc\_ms} field exists in the JSON output but is always $0.0$---GPU union-find was planned but not implemented.  The header \texttt{tile\_kernel.cuh} is named suggestively but contains a primality kernel from an earlier iteration, not a union-find kernel.

Moving union-find to GPU (using ECL-CC~\cite{jaiganesh2018} or Afforest~\cite{sutton2018} algorithms) eliminates the device-to-host bitmap copy and the CPU serialization bottleneck.  The tile grid ($2000 \times 2000$, ${\sim}\,4\text{M}$ cells) fits comfortably in GPU shared memory or L2 cache, and the connectivity pattern (128 neighbors for $k^2 = 40$) produces a sparse but regular graph.

\subsubsection{Tile-per-SM batching}

A single $2000 \times 2000$ tile generates ${\sim}\,2{,}013$ CUDA blocks.  On a 4090 with 128~SMs, this underutilizes the GPU (only ${\sim}\,16$ blocks/SM).  With a fully GPU-side pipeline, multiple tiles can be processed concurrently---one tile per SM---achieving full hardware utilization.

\subsubsection{Throughput projections}

\begin{table}[ht]
\centering
\begin{tabular}{@{}lccrr@{}}
\toprule
Hardware & SMs & Clock (MHz) & Per-tile (batched) & 589K tiles \\
\midrule
Jetson Orin Nano & 16  & 918  & ${\sim}\,180\text{ ms}$ & ${\sim}\,29\text{ hr}$ \\
RTX 4090         & 128 & 2520 & ${\sim}\,1.8\text{ ms}$ & ${\sim}\,18\text{ min}$ \\
A100 SXM4        & 108 & 1410 & ${\sim}\,2.2\text{ ms}$ & ${\sim}\,22\text{ min}$ \\
H100 SXM5        & 132 & 1980 & ${\sim}\,1.5\text{ ms}$ & ${\sim}\,15\text{ min}$ \\
\bottomrule
\end{tabular}
\caption{Projected throughput with fully-GPU tile pipeline and tile-per-SM batching.  The 589K tiles column shows the wall time for a single-shell UB at $R = 1.5\text{B}$ (first-octant coverage).}
\label{tab:throughput-projection}
\end{table}

\subsubsection{UB campaign cost}

\begin{table}[ht]
\centering
\begin{tabular}{@{}lrrrr@{}}
\toprule
Scenario & Tiles & $1\times$ 4090 & $8\times$ 4090 & $8\times$ A100 \\
\midrule
Single shell ($R = 1.5\text{B}$) & 589K & 18 min & 2.3 min & 2.8 min \\
1M radial (worst case) & 294.5M & 8.5 days / \$408 & 25 hr / \$400 & 11 hr / \$176 \\
\bottomrule
\end{tabular}
\caption{UB campaign cost projections with GPU union-find and tile-per-SM batching.}
\label{tab:ub-campaign-cost}
\end{table}

\subsubsection{CPU bottleneck on large GPUs}

On fast GPUs, the CPU union-find becomes the bottleneck even with batching: at $40\text{ ms/tile} \times 589\text{K tiles} / n_{\text{cores}}$, a 64-core CPU still serializes to ${\sim}\,6.1\text{ min}$ for the UF step alone---comparable to the GPU primality time on a 4090.  A fully GPU-side pipeline eliminates this entirely; no bitmap leaves the device until the final face-port summary.

\subsubsection{Implementation status}

GPU union-find for grid graphs is a bounded engineering task using well-studied algorithms (ECL-CC, Afforest).  The gap is documented; the method's computational feasibility does not depend on novel GPU research.  The \texttt{gpu\_cc\_ms} field provides a natural integration point for timing instrumentation once the kernel is implemented.

\subsection{Status}

The UB method is formalized but not yet deployed as a campaign.  The existing \texttt{compose.rs} implements distance-based matching (sound for UB).  Shared-prime matching, which would yield tighter composition, is specified but not yet implemented.  Both adversarial reviews (Sections on tile-UB and ISE challenges) returned verdicts of ``Accept with revisions.''


% ══════════════════════════════════════════════════════════════
\section{Fat Stripe: Tile-Composed Annular Moat Detection}\label{sec:fat-stripe}

The fat-stripe method is a concrete realization of the tile-based UB approach (Section~\ref{sec:ub}).  Rather than leaving gapless tiling as a feasibility sketch, fat stripe implements it end-to-end: tile the first-octant annular strip, process each tile independently, compose all tiles via the GPCTO algebra, and read off a definitive verdict on whether any connected component of~$\Gk$ spans the annulus radially.  Unlike the ISE, which samples sparse angular positions and measures local percolation, fat stripe covers the full first-octant angular extent and answers a strictly stronger question: \emph{does any component---origin-connected or not---bridge the annular band?}

\subsection{Annular strip geometry}

\begin{definition}[Annular strip]\label{def:annular-strip}
Fix radii $r_{\min} < r_{\max}$ and step bound~$k^2$.  The \emph{first-octant annular strip} is the lattice region
\[
  \mathcal{A} = \{a + bi \in \Zi : r_{\min} \le a \le r_{\max},\; 0 \le b \le b_{\max}(a)\},
\]
where $b_{\max}(a) = \min\!\big(\lceil\sqrt{r_{\max}^2 - a^2}\rceil,\; a\big)$ enforces both the circular outer boundary and the first-octant constraint $b \le a$.
\end{definition}

\begin{remark}[Angular span]\label{rem:angular-span}
The angular half-width of~$\mathcal{A}$ at radius~$R$ is approximately
\[
  \theta \approx \arctan\!\left(\frac{\sqrt{2R \cdot \Delta r}}{R}\right) \approx \sqrt{\frac{2\Delta r}{R}},
\]
where $\Delta r = r_{\max} - r_{\min}$.  For $\Delta r / R \sim 10^{-4}$ (a 10K-wide annulus at $R = 80\text{M}$), $\theta \approx 0.8^\circ$, covering roughly $11\%$ of the first octant.  The 8-fold exact symmetry of~$\Zi$ under the dihedral group~$D_4$ (Lemma~\ref{lem:octant}) ensures that a barrier in this sector implies a barrier in all eight octants, provided the full octant angular range is covered; see Section~\ref{sec:fs-geometry} for a discussion of partial-octant coverage.
\end{remark}

\subsection{Tiling and tile processing}

The annular strip~$\mathcal{A}$ is partitioned into a grid of tiles.  The radial dimension is divided into \emph{stripes} of height~$H$ (typically $H = W$, the tile width), and each stripe is divided laterally into tiles of width~$W$.

\begin{definition}[Radial stripe]\label{def:radial-stripe}
A \emph{radial stripe} is the sub-band
\[
  \mathcal{A}_n = \{a + bi \in \mathcal{A} : a_n \le a < a_n + H\},
\]
where $a_n = r_{\min} + nH$ for $n = 0, 1, \ldots, N_r - 1$ and $N_r = \lceil(r_{\max} - r_{\min})/H\rceil$.
\end{definition}

Within each stripe, tiles are indexed laterally: $T_{n,j} = T(a_n, b_j, W, H)$ with $b_j = jW$, for $j = 0, 1, \ldots, J_n - 1$ where $J_n = \lceil b_{\max}(a_n) / W \rceil$.

Each tile is processed independently in three phases:
\begin{enumerate}
  \item \textbf{Row sieve + Miller--Rabin.}  For each row (fixed~$a$) of the expanded tile~$T_{n,j}^+$, a sieve with splitting primes up to $L = 110{,}000$ eliminates composite Gaussian norms, and deterministic Miller--Rabin (4--12 witnesses depending on norm magnitude) confirms surviving candidates.  This produces a prime bitmap for the expanded tile.

  \item \textbf{Sparse union-find.}  Primes are compacted from the bitmap into a list of approximately $P \approx 4S^2 / (\pi \ln n)$ entries (where $S = W + 2c$ is the expanded tile side and~$n$ is the typical norm).  A rank table constructed via block-wise popcount enables $O(1)$ neighbor-index lookup.  Union-find over the compacted list---checking all backward offset vectors $(\Delta a, \Delta b)$ with $\Delta a^2 + \Delta b^2 \le k^2$---identifies connected components.  Working set: approximately $500\text{ KB}$ (rank table + UF arrays), versus $16\text{ MB}$ for a dense UF over all lattice points.

  \item \textbf{Face-port extraction.}  For each connected component, the algorithm records which faces ($\fI$, $\fO$, $\fL$, $\fR$) it touches within the interior (non-collar) tile, together with the $(a, b)$ coordinates of the face-port primes and their component identifiers.  This produces a \texttt{TileOperator} as defined in Section~\ref{sec:tile-operator}.
\end{enumerate}

\subsection{Hierarchical composition}

Tile operators are composed in a two-level hierarchy that mirrors the strip-and-tile geometry:

\paragraph{Step 1: Horizontal composition within stripes.}
Within each radial stripe~$\mathcal{A}_n$, tiles are composed left-to-right via $\otimes_h$ (Definition~\ref{def:h-composition}).  For efficiency, tiles are grouped into \emph{column-chunks} of~$C$ consecutive tiles (typically $C = 200$); tiles within each chunk are composed sequentially, then inter-chunk seams are merged.  The result is a single band operator~$B_n$ for each radial stripe.

\paragraph{Step 2: Vertical composition across stripes.}
Band operators are composed bottom-to-top via $\otimes_v$ (Definition~\ref{def:v-composition}), producing a single operator~$B$ for the full annular strip.

\begin{theorem}[Composition order independence]\label{thm:fs-order}
The final composed operator~$B$ is independent of the order in which horizontal and vertical compositions are evaluated.
\end{theorem}

\begin{proof}
By Theorem~\ref{thm:composition-sound}, each composition step merges components via union-find on face-port distances.  Union-find is commutative and associative: the final partition depends only on the set of union operations, not their order.  Any evaluation schedule---left-to-right within rows then bottom-to-top across rows, binary reduction trees, or any hybrid---produces the same set of face-port pairs satisfying $|p_1 - p_2|^2 \le k^2$, hence the same final partition.
\end{proof}

\subsection{Spanning verdict}

\begin{definition}[Radial spanning]\label{def:radial-spanning}
A connected component~$C$ of the composed operator~$B$ is \emph{radially spanning} if there exist face ports $p_{\text{in}}, p_{\text{out}} \in C$ such that
\[
  |p_{\text{in}}| \le r_{\min} + c \qquad \text{and} \qquad |p_{\text{out}}| \ge r_{\max} - c,
\]
where $|\cdot|$ denotes the Euclidean norm $\sqrt{a^2 + b^2}$ and~$c$ is the collar width.
\end{definition}

\begin{remark}[Radius-based vs.\ face-based verdict]\label{rem:radius-verdict}
Because the tiling is Cartesian while the moat is circular, the Inner and Outer face labels of the composed operator do not correspond exactly to the circular boundaries $|z| = r_{\min}$ and $|z| = r_{\max}$.  The spanning verdict therefore inspects the actual $(a, b)$ coordinates of all face ports, computing $\sqrt{a^2 + b^2}$ for each, rather than relying on the rectangular face assignments.  This avoids false verdicts near the corners of the tiled region where the circular and rectangular boundaries diverge.
\end{remark}

\begin{theorem}[Soundness of the spanning verdict]\label{thm:fs-soundness}
If no component of~$B$ is radially spanning (Definition~\ref{def:radial-spanning}), then no connected component of~$\Gk$ restricted to the first-octant annular strip~$\mathcal{A}$ has a path from $|z| \le r_{\min}$ to $|z| \ge r_{\max}$.
\end{theorem}

\begin{proof}
Suppose for contradiction that a path $\pi_1, \pi_2, \ldots, \pi_m$ in $\Gk[\mathcal{A}]$ connects some~$\pi_1$ with $|\pi_1| \le r_{\min}$ to~$\pi_m$ with $|\pi_m| \ge r_{\max}$.  Every edge $\{\pi_i, \pi_{i+1}\}$ satisfies $|\pi_i - \pi_{i+1}|^2 \le k^2$ and both endpoints lie in~$\mathcal{A}$.

By the gapless tiling (Definition~\ref{def:gapless}), every prime in~$\mathcal{A}$ belongs to the interior of at least one tile~$T_{n,j}$.  Each consecutive pair $(\pi_i, \pi_{i+1})$ either lies in the same expanded tile (captured by the local UF) or straddles a tile boundary (captured by composition, since both endpoints are face ports in their respective tiles).  Therefore, all~$\pi_i$ lie in the same component of~$B$.

The first prime~$\pi_1$ has $|\pi_1| \le r_{\min}$, making it a face port satisfying $|p| \le r_{\min} + c$ (it lies in the face region of the lowest stripe).  The last prime~$\pi_m$ has $|\pi_m| \ge r_{\max}$, making it a face port satisfying $|p| \ge r_{\max} - c$.  Hence their shared component is radially spanning---contradicting the assumption.
\end{proof}

\begin{corollary}[Conservative direction]\label{cor:fs-conservative}
The fat-stripe verdict ``blocked'' (no spanning component) is conservative: it can only be declared when no component of~$\Gk$ crosses the annulus within the tiled region.  False ``blocked'' verdicts are impossible; false ``not blocked'' verdicts can occur only due to missed connections from distance-based composition (Remark~\ref{rem:conservative}).
\end{corollary}

\subsection{Moat types and the percolation regime}

The ISE (Section~\ref{sec:ise}) and fat stripe answer related but distinct questions.  The ISE measures whether \emph{any} component crosses a single tile's radial extent; fat stripe measures whether \emph{any} component spans an entire annular strip.  These yield different detection capabilities depending on the percolation regime of~$\Gk$ at the candidate moat radius.

\begin{definition}[Total barrier]\label{def:total-barrier}
A moat at radius~$R$ is a \emph{total barrier} if no connected component of~$\Gk$---whether origin-connected or not---has a path from $|z| < R$ to $|z| > R + \Delta r$, for some annular width $\Delta r > 0$.
\end{definition}

\begin{definition}[Origin-only moat]\label{def:origin-only}
A moat at radius~$R$ is \emph{origin-only} if the origin component $C_0(k)$ does not reach $|z| > R$, but other connected components of~$\Gk$ do have paths crossing the annular region around~$R$.
\end{definition}

Whether a moat is a total barrier or origin-only depends on the local degree statistics.  Degree values for $k^2 \ge 32$ were measured empirically using the \texttt{--degree-stats} tool in \texttt{fat-stripe}, which samples backward and forward neighbor counts at each Gaussian prime within the annulus.  Values for $k^2 \le 26$ are estimated from $\rho \times |\text{offsets}|$.

\begin{table}[ht]
\centering
\small
\begin{tabular}{@{}cccccccl@{}}
\toprule
$k^2$ & $R_{\text{moat}}$ & Bkwd offsets & Mean bkwd deg & Mean total deg & Isolated \% & Moat type \\
\midrule
2  & ${\sim}\,12$       & 2  & 0.45 & 0.90 & --- & Total barrier \\
4  & ${\sim}\,45$       & 2  & 0.31 & 0.61 & --- & Total barrier \\
26 & ${\sim}\,1.02\text{M}$ & 40 & ${\sim}\,1.0$ & ${\sim}\,2.0$ & --- & Origin-only \\
32 & ${\sim}\,2.82\text{M}$ & 50 & 1.96 & 3.92 & 1.3\% & Total barrier \\
36 & ${\sim}\,80\text{M}$   & 56 & 1.99 & 3.97 & 1.3\% & Total barrier \\
40 & ${\sim}\,1.05\text{B}$ & 64 & 2.02 & 4.03 & 1.3\% & Total barrier (confirmed) \\
\bottomrule
\end{tabular}
\caption{Degree statistics and moat type classification at known and confirmed moat radii.  Degrees for $k^2 \ge 32$ are empirically measured via \texttt{--degree-stats}.}
\label{tab:moat-types}
\end{table}

Empirical degree measurements reveal a universal critical threshold: the percolation transition (total annular blockage) occurs at $d_c \approx 4.0$ mean total degree across all measured $k^2$~values.  At confirmed moat radii, the mean total degree sits just below~$4.0$ ($3.92$--$3.97$ for $k^2 = 32, 36$); at supercritical radii, it exceeds~$4.0$ ($4.10$--$4.18$).

At $k^2 = 26$, the total degree at the moat radius is approximately~$2.0$, placing the graph in the regime where the origin component is bounded but orphan clusters still span the annulus.  Fat stripe reports ``not blocked'' for such moats, correctly reflecting that the annulus is not a total barrier.

At $k^2 = 32$ and $k^2 = 36$, the measured total degree at the moat radius is $3.92$--$3.97$, just below the $d_c \approx 4.0$ threshold.  The graph is sufficiently sparse that large-scale fragmentation prevents \emph{any} component from spanning the annulus, not merely the origin component.  Fat stripe detects these as total barriers.

\begin{remark}[Complementarity with ISE]\label{rem:complementarity}
For origin-only moats ($k^2 = 26$), the ISE detects $\fr$ depressions but not $\fr = 0$, because orphan clusters cross the annulus.  Fat stripe confirms ``not blocked.''  For total-barrier moats ($k^2 = 32, 36$), both methods agree: ISE shows the percolation transition, and fat stripe confirms absolute disconnection.  The two methods are complementary: ISE scans efficiently for moat candidates across wide radial ranges; fat stripe provides the definitive verdict at specific candidate radii.
\end{remark}

\subsection{Calibration results}\label{sec:fs-calibration}

Fat-stripe calibration runs were performed on a Mac Mini (M-series) using the \texttt{fat-stripe} crate with $L = 110{,}000$ sieve limit, deterministic Miller--Rabin, and sparse union-find.

\subsubsection{\texorpdfstring{$k^2 = 32$}{k\texttwosuperior\ = 32} at \texorpdfstring{$R \approx 2.82\text{M}$}{R approx 2.82M}}

The known Tsuchimura lower bound for $k^2 = 32$ is $R_{\text{moat}} \approx 2{,}823{,}054$ (the origin component reaches this radius but no further).

\begin{table}[ht]
\centering
\begin{tabular}{@{}lr@{}}
\toprule
Parameter & Value \\
\midrule
Annulus                & $[2{,}820{,}000,\; 2{,}830{,}000]$ (width $10{,}000$) \\
Tile size              & $2000 \times 2000$ \\
Radial stripes         & 5 \\
Total tiles            & 449 \\
Column-chunk size      & 200 \\
Angular span           & $4.82^\circ$ (${\approx}\,11\%$ of first octant) \\
Elapsed time           & 34 seconds \\
\midrule
\textbf{Verdict}       & \textbf{BLOCKED --- 0 spanning components} \\
\bottomrule
\end{tabular}
\caption{Fat-stripe calibration at $k^2 = 32$, $R \approx 2.82\text{M}$.}
\label{tab:fs-k32}
\end{table}

\subsubsection{\texorpdfstring{$k^2 = 36$}{k\texttwosuperior\ = 36} at \texorpdfstring{$R \approx 80\text{M}$}{R approx 80M}}

The known Tsuchimura upper bound for $k^2 = 36$ is $R_{\text{moat}} < 80{,}015{,}782$.

\begin{table}[ht]
\centering
\begin{tabular}{@{}lr@{}}
\toprule
Parameter & Value \\
\midrule
Annulus                & $[80{,}005{,}000,\; 80{,}015{,}000]$ (width $10{,}000$) \\
Tile size              & $2000 \times 2000$ \\
Radial stripes         & 5 \\
Total tiles            & $2{,}373$ \\
Column-chunk size      & 200 \\
Elapsed time           & 4 minutes 5 seconds \\
\midrule
\textbf{Verdict}       & \textbf{BLOCKED --- 0 spanning components} \\
\bottomrule
\end{tabular}
\caption{Fat-stripe calibration at $k^2 = 36$, $R \approx 80\text{M}$.}
\label{tab:fs-k36}
\end{table}

Both calibration runs confirm that at the known moat radii for $k^2 = 32$ and $k^2 = 36$, the annular strip is a total barrier: no connected component of~$\Gk$ spans from the inner to the outer radius within the tiled first-octant region.

\subsubsection{\texorpdfstring{$k^2 = 40$}{k\texttwosuperior\ = 40} at \texorpdfstring{$R \in [1.03\text{B},\; 1.09\text{B}]$}{R in [1.03B, 1.09B]}}\label{sec:fs-k40}

No moat at $k^2 = 40$ was previously known.  The ISE sigmoid (Section~\ref{sec:k40-results}) with $\Rhalf \approx 839\text{M}$ and the percolation parameter analysis (Section~\ref{sec:fs-k40-implications}) predicted that the candidate moat radius lies near $R \approx 1.05\text{B}$ in the same fragmentation regime as $k^2 = 36$.  Five fat-stripe probes were run across a 60M radial band to test this prediction.

Each probe covers a $128{,}000 \times 128{,}000$ annular strip ($64 \times 64$ tiles at $W = 2000$), with angular width $b_{\max} = 128{,}000$.

\begin{table}[ht]
\centering
\begin{tabular}{@{}ccccc@{}}
\toprule
Probe & $R$ & Spanning components & Tiles & Time (Mac) \\
\midrule
1 & $1.03\text{B}$ & 0 & $4{,}096$ & 5m\,15s \\
2 & $1.04\text{B}$ & 0 & $4{,}096$ & 5m\,15s \\
3 & $1.05\text{B}$ & 0 & $4{,}096$ & 5m\,13s \\
4 & $1.07\text{B}$ & 0 & $4{,}096$ & 5m\,12s \\
5 & $1.09\text{B}$ & 0 & $4{,}096$ & 5m\,9s \\
\midrule
\multicolumn{2}{@{}l}{\textbf{Verdict}} & \multicolumn{3}{l}{\textbf{BLOCKED --- 0 spanning components (all probes)}} \\
\bottomrule
\end{tabular}
\caption{Fat-stripe probes at $k^2 = 40$, $R \in [1.03\text{B}, 1.09\text{B}]$.}
\label{tab:fs-k40}
\end{table}

Per-tile throughput: ${\sim}\,0.076\text{ s/tile}$ on Mac (M-series), comparable to the $k^2 = 36$ calibration run.

\paragraph{Angular coverage.}
At $R = 1\text{B}$, the 128K angular width subtends ${\sim}\,0.007^\circ$, or ${\sim}\,0.016\%$ of the first octant.  This is a narrow-wedge probe---sufficient to detect the fragmentation regime and confirm methodological consistency with $k^2 = 32$ and $k^2 = 36$, but far below the full-octant coverage (Lemma~\ref{lem:fs-octant}) required for a rigorous moat proof.

\paragraph{Significance.}
This is the first detection of total annular blockage at $k^2 = 40$.  Unlike $k^2 = 32$ and $k^2 = 36$, where the moat was independently known from Tsuchimura's exhaustive method, $k^2 = 40$ has no prior moat result.  The fat-stripe probes advance from calibration (confirming known results) to discovery (detecting a new candidate moat region).

\subsubsection{Cross-validation with ISE at \texorpdfstring{$k^2 = 32$}{k\texttwosuperior\ = 32}}

ISE probes with $200 \times 200$ tiles and 32 strips at the same radial location show per-tile $\fr$ values ranging from $0.31$ to $0.78$---individual small tiles have $31$--$78\%$ inner-to-outer connectivity.  This is \emph{not} contradictory with the fat-stripe ``blocked'' verdict.  Individual tiles are well-connected \emph{locally} (most components within a single tile span its 200-unit radial extent), but \emph{long-range} radial connectivity through the full 10,000-unit composed strip breaks down at the moat distance.  The moat is a large-scale fragmentation phenomenon, not a local one: components that bridge a single tile's height become disconnected from one another when the full angular extent is considered.

\subsection{Geometry and symmetry}\label{sec:fs-geometry}

\paragraph{Partial-octant coverage.}
The angular span of the fat-stripe annulus at $R = r$ with width $\Delta r$ is approximately $\theta \approx \sqrt{2\Delta r / R}$, which for the calibration runs amounts to $5$--$11\%$ of the first octant.  The full octant spans~$45^\circ$, so the calibrated region covers roughly $1$--$5^\circ$ out of~$45^\circ$.

\begin{lemma}[Octant symmetry for total barriers]\label{lem:fs-octant}
Let~$\mathcal{A}_\theta$ denote the first-octant annular strip restricted to angular positions $[0, \theta]$.  If no component of~$\Gk$ restricted to~$\mathcal{A}_\theta$ spans from $r_{\min}$ to $r_{\max}$, this does not immediately imply a full-annulus barrier.  However, by the 8-fold dihedral symmetry of~$\Zi$, an identical argument applies to each of the 8 octant copies of~$\mathcal{A}_\theta$.  A total annular barrier is proven if and only if: \textup{(a)}~the tiled region covers the full first-octant angular extent ($0 \le b \le a$), or \textup{(b)}~separate fat-stripe runs cover complementary angular sectors whose union tiles the full octant gaplessly.
\end{lemma}

\begin{remark}[Calibration vs.\ proof]\label{rem:calibration-vs-proof}
The calibration results of Section~\ref{sec:fs-calibration} demonstrate that fat stripe correctly identifies total barriers at known moat radii ($k^2 = 32, 36$) and detects a new candidate barrier ($k^2 = 40$), but all probes tile only a sub-sector of the first octant.  A rigorous moat proof requires full-octant coverage (Lemma~\ref{lem:fs-octant}).  For $k^2 = 32$ and $k^2 = 36$, the sub-sector verdicts confirm methodological correctness against independently verified Tsuchimura moats.  For $k^2 = 40$, the narrow-wedge probes (${\sim}\,0.007^\circ$) demonstrate that total blockage occurs in the predicted fragmentation regime, but full-octant tiling is required for a rigorous moat claim.
\end{remark}

\subsection{Implications for \texorpdfstring{$k^2 = 40$}{k\texttwosuperior\ = 40}}\label{sec:fs-k40-implications}

Empirical degree statistics measured via \texttt{fat-stripe --degree-stats} at the moat radii for $k^2 = 32$, $36$, and $40$ reveal a universal critical degree threshold:

\begin{table}[ht]
\centering
\begin{tabular}{@{}lccc@{}}
\toprule
Parameter & $k^2 = 32$ at $R = 2.82\text{M}$ & $k^2 = 36$ at $R = 80\text{M}$ & $k^2 = 40$ at $R = 1.05\text{B}$ \\
\midrule
Backward offsets            & 50   & 56   & 64   \\
Mean backward degree        & 1.96 & 1.99 & 2.02 \\
Mean total degree           & 3.92 & 3.97 & 4.03 \\
Isolated primes             & 1.3\% & 1.3\% & 1.3\% \\
\bottomrule
\end{tabular}
\caption{Empirical degree statistics at moat radii, measured via \texttt{--degree-stats}.}
\label{tab:k40-comparison}
\end{table}

All three $k^2$~values show mean total degree $\approx 4.0$ at or near the moat radius, confirming $d_c \approx 4.0$ as the universal critical threshold.  At supercritical radii (well inside the moat), degrees rise to $4.10$--$4.18$; at confirmed moat radii, they sit at $3.92$--$3.97$.  The transition is sharp: $k^2 = 40$ crosses $d_c \approx 4.0$ in the range $R \approx 1.0\text{B}$--$1.2\text{B}$, consistent with both the ISE sigmoid ($\Rhalf \approx 839\text{M}$) and the fat-stripe probes of Section~\ref{sec:fs-k40} showing zero spanning components across $R \in [1.03\text{B}, 1.09\text{B}]$.

\paragraph{Scaling considerations.}
A full-octant fat-stripe campaign at $R \approx 1.05\text{B}$ with $\Delta r = 10{,}000$ requires approximately $J = \lceil 1.05 \times 10^9 / 2000 \rceil = 525{,}000$ tiles per stripe and $N_r = 5$ radial stripes, totaling ${\sim}\,2.6\text{M}$~tiles.  At the calibrated throughput of ${\sim}\,0.5\text{ s/tile}$ (Mac Mini), this amounts to approximately $360$~hours serial.  With Rayon parallelism across 8--64 cores and a double-buffered pipeline, wall-clock time reduces to 6--45 hours depending on hardware.  Cloud deployment on multi-core instances makes a single-day campaign feasible.


% ══════════════════════════════════════════════════════════════
\section{Empirical Validation}\label{sec:validation}

\subsection{Calibration against known moats}

ISE calibration runs were performed on a Jetson Orin Nano with $M = 32$ positive-only strips, square tiles, and the scanline kernel (Montgomery multiplication, row-sieve, union-find).

\begin{table}[ht]
\centering
\small
\begin{tabular}{@{}cccccccc@{}}
\toprule
$k^2$ & Tile size & Radial range & Shells & Known moat $R$ & ISE $\fr$ at moat & Control $\bar{f}$ & Moat sweep $\bar{f}$ \\
\midrule
2  & $8 \times 8$       & $[0, 100]$                     & 13  & ${\sim}\,12$      & $0.000$ at $R{=}20$       & ---    & $0.19$ \\
26 & $500 \times 500$   & $[950\text{K}, 1.1\text{M}]$   & 300 & $1{,}015{,}639$   & $0.250$ at $R{\approx}1{,}016{,}250$ & $0.937$ & $0.477$ \\
32 & $500 \times 500$   & $[2.7\text{M}, 2.9\text{M}]$   & 400 & $2{,}823{,}055$   & $0.344$ at $R{\approx}2{,}823{,}250$ & $0.891$ & $0.503$ \\
36 & $2000 \times 2000$ & $[79.5\text{M}, 80.5\text{M}]$ & 500 & Unknown           & $\mathbf{0.000}$ at $R{=}80{,}399{,}000$ & $0.881$ & $0.212$ \\
\bottomrule
\end{tabular}
\caption{ISE calibration against known Tsuchimura moats.}
\label{tab:calibration}
\end{table}

Key findings:
\begin{enumerate}
  \item \textbf{$k^2 = 2$:}  ISE correctly detects the moat with six $\fr = 0$ shells spanning $R = 20$--$76$.
  \item \textbf{$k^2 = 26, 32$:}  ISE detects $\fr$ depressions near the Tsuchimura moats but does not reach $\fr = 0$.  This is expected: the moat is a structural gap in the origin component, not a collapse of all local connectivity.  Orphan clusters provide local IO crossings even at the moat radius.  The ISE correctly distinguishes control regions ($\bar{f} > 0.88$) from moat regions ($\bar{f} \approx 0.50$).
  \item \textbf{$k^2 = 36$:}  A single $\fr = 0$ shell at $R = 80{,}399{,}000$ with monotonic approach ($f = 0.156 \to 0.125 \to 0.063 \to 0.000$).  This marks the percolation transition boundary.
\end{enumerate}

\subsection{Tile-size sensitivity}

Four tile heights at $k^2 = 26$, $R \in [1\text{M}, 1.03\text{M}]$, 32 strips:

\begin{table}[ht]
\centering
\begin{tabular}{@{}ccccc@{}}
\toprule
Tile & Shells & $f_{\min}$ & $f_{\max}$ & Time/shell \\
\midrule
$200 \times 200$   & 150 & 0.281 & 0.750 & 0.7s \\
$500 \times 500$   & 60  & 0.250 & 0.625 & 3.0s \\
$1000 \times 1000$ & 30  & 0.313 & 0.594 & 11.1s \\
$2000 \times 2000$ & 15  & 0.313 & 0.594 & 41.9s \\
\bottomrule
\end{tabular}
\caption{Tile-size sensitivity at $k^2 = 26$.}
\label{tab:tile-sensitivity}
\end{table}

Tile height does not significantly shift~$\fr$ estimates.  The mean and extrema are consistent across two orders of magnitude in tile area.  Larger tiles reduce variance (more primes per tile) but do not shift the mean.

\subsection{\texorpdfstring{$k^2 = 40$}{k\texttwosuperior\ = 40} results}\label{sec:k40-results}

ISE probes at $k^2 = 40$ with $M = 64$ strips, $2000 \times 2000$ tiles, from $R = 600\text{M}$ to $R = 3.5\text{B}$:

\begin{table}[ht]
\centering
\begin{tabular}{@{}ccl@{}}
\toprule
$R$ & $\fr$ & Regime \\
\midrule
$600\text{M}$ & $1.000$ & Connected \\
$700\text{M}$ & $0.906$ & Declining \\
$800\text{M}$ & $0.609$ & Transition \\
$900\text{M}$ & $0.391$ & Transition \\
$1.0\text{B}$ & $0.203$ & Near collapse \\
$1.2\text{B}$ & $0.078$ & Sparse crossings \\
$1.5\text{B}$ & $0.016$ & Near extinction \\
$2.0\text{B}$ & $0.000$ & Extinct \\
$2.5\text{B}$ & $0.000$ & Extinct \\
$3.0\text{B}$ & $0.000$ & Extinct \\
$3.5\text{B}$ & $0.000$ & Extinct \\
\bottomrule
\end{tabular}
\caption{ISE $\fr$ profile for $k^2 = 40$.}
\label{tab:k40-results}
\end{table}

The $\fr$ profile is a sigmoid with:
\begin{itemize}
  \item $\Rhalf \approx 839\text{M} \pm 12\text{M}$ (95\% CI from logistic fit)
  \item First $\fr = 0$ shells at $R \approx 2.0\text{B}$
  \item Total extinction from $2.0\text{B}$ onward
\end{itemize}

The logistic fit uses three calibration points in the transition region ($R = 700\text{M}, 800\text{M}, 900\text{M}$).  Alternative functional forms (complementary error function, asymmetric sigmoid) have not yet been compared; model-selection sensitivity is an open question.


% ══════════════════════════════════════════════════════════════
\section{Open Questions}\label{sec:open}

\begin{enumerate}
  \item \textbf{Finite-size scaling.}  Does $\Rhalf$ depend on tile size~$W$?  Data at $W \in \{500, 1000, 2000, 4000\}$ exists for $k^2 = 26$ but formal extrapolation to $W \to \infty$ has not been performed.  The ISE moat estimate~$\Rhalf$ should be understood as a property of a specific tile geometry, not intrinsically of~$\Gk$.  A finite-size scaling protocol---measuring $\Rhalf(W)$ for several widths and extrapolating---would yield a tile-size-independent quantity.

  \item \textbf{Percolation theory connection.}  Is the tile connectivity transfer operator related to existing operators in percolation theory (e.g., the transfer matrix for site percolation on~$\mathbb{Z}^2$)?  The analogy is suggestive---both compress connectivity information through a strip---but Gaussian primes are deterministic, not random, so the connection is heuristic.

  \item \textbf{Optimal tile geometry.}  Square tiles ($W = H$) are confirmed necessary for reliable IO crossing detection.  Is there an optimal~$W$ for a given~$k^2$ and radius~$R$?  The correlation length~$\xi$ of the crossing process sets a natural scale: $W \gg \xi$ is needed for reliable detection, but $W \gg \xi$ wastes computation.

  \item \textbf{GPU-parallel UB.}  Can the UB method be made GPU-parallel with composition on device?  The individual tile builds are embarrassingly parallel (one CUDA thread block per tile), but the horizontal composition step requires sequential union-find merging across~$N_\ell$ tiles.  A hierarchical super-tile aggregation could reduce this to $O(\log N_\ell)$ sequential composition steps.

  \item \textbf{Logistic model validation.}  The sigmoid fit to~$\fr$ data is based on a logistic function fitted to a narrow transition window.  The complementary error function $\tfrac{1}{2}\operatorname{erfc}(\cdot)$ is arguably more physically motivated for a percolation-like transition.  Fitting both models and comparing~$\Rhalf$ estimates via AIC/BIC would quantify model-choice uncertainty.

  \item \textbf{Tile-size dependence of $\Rhalf$.}  Narrower tiles measure lower~$\fr$ (they miss lateral paths), shifting the sigmoid leftward.  If $\Rhalf$ at $W = 2000$ and $\Rhalf$ at $W = 4000$ differ by more than the CI, the estimation framework requires finite-size scaling corrections.  This is the deepest methodological concern.
\end{enumerate}


% ══════════════════════════════════════════════════════════════
\appendix

\section{Notation Summary}\label{app:notation}

\begin{table}[ht]
\centering
\begin{tabular}{@{}ll@{}}
\toprule
Symbol & Definition \\
\midrule
$\Zi$                            & Gaussian integers \\
$P$                              & Set of all Gaussian primes \\
$\Gk = (P, E_k)$                & Gaussian prime graph at step bound $k^2$ \\
$\Sigma_k$                       & Seed set: primes with $|\pi|^2 \le k^2$ \\
$C_0(k)$                         & Origin component of $\Gk$ \\
$\Rmoat(k)$                      & Moat radius: supremum of radii reached by $C_0(k)$ \\
$T(a_0, b_0, W, H)$             & Tile: lattice rectangle $[a_0, a_0{+}H] \times [b_0, b_0{+}W]$ \\
$c = \lceil\sqrt{k^2}\rceil$    & Collar width \\
$\Tplus$                         & Expanded tile (interior + collar of width $c$) \\
$\fI, \fO, \fL, \fR$            & Inner, Outer, Left, Right faces \\
$\spant(X,Y)$                    & Count of components touching both faces $X$ and $Y$ \\
$\iocount(T)$                    & $= \spant(\fI, \fO)$ \\
$S_j$                            & Strip $j$: vertical stack of tiles at lateral offset $b_j$ \\
$M$                              & Number of strips \\
$\fr$                            & Fraction of strips with $\iocount > 0$ at radius $r$ \\
$\Rhalf$                         & Radius where $\fr = 0.5$ (logistic midpoint) \\
$\rho(R)$                        & Local Gaussian prime density at radius $R$ \\
$\mathcal{A}$                    & First-octant annular strip \\
$\mathcal{A}_n$                  & Radial stripe $n$ within the annular strip \\
$B_n$                            & Band operator: horizontally composed tiles within stripe $n$ \\
$B$                              & Full composed operator for the annular strip \\
\bottomrule
\end{tabular}
\caption{Notation summary.}
\label{tab:notation}
\end{table}


% ══════════════════════════════════════════════════════════════
\section{Algorithm Pseudocode (ISE)}\label{app:pseudocode}

\subsection{Tile operator computation}

\begin{algorithm}[H]
\caption{ComputeTile$(a_0, b_0, W, H, k^2)$}\label{alg:compute-tile}
\KwIn{Tile bounds $(a_0, b_0, W, H)$, step bound $k^2$}
\KwOut{TileOperator (component face sets, face ports, $\iocount$)}
$c \gets \lceil\sqrt{k^2}\rceil$\;
Enumerate all Gaussian primes $\pi$ in the expanded rectangle $[a_0 - c,\; a_0 + H + c] \times [b_0 - c,\; b_0 + W + c]$\;
Initialize union-find \texttt{UF} on $|\text{primes}|$ elements\;
Build spatial hash: for each $\pi$, insert into cell grid with cell size $c$\;
\ForEach{prime $\pi$ at $(a, b)$}{
  \ForEach{neighbor cell within distance 2 cells}{
    \ForEach{prime $\pi'$ in that cell with $\mathrm{index}(\pi') > \mathrm{index}(\pi)$}{
      \If{$|\pi - \pi'|^2 \le k^2$}{
        $\texttt{UF.Union}(\mathrm{index}(\pi),\, \mathrm{index}(\pi'))$\;
      }
    }
  }
}
\ForEach{prime $\pi = (a, b)$ in the interior tile $[a_0, a_0{+}H] \times [b_0, b_0{+}W]$}{
  $\text{root} \gets \texttt{UF.Find}(\mathrm{index}(\pi))$\;
  $\text{component} \gets \text{ComponentID}(\text{root})$\;
  \lIf{$a - a_0 \le c$}{tag component with $\fI$; add face port}
  \lIf{$(a_0 + H) - a \le c$}{tag component with $\fO$; add face port}
  \lIf{$b - b_0 \le c$}{tag component with $\fL$; add face port}
  \lIf{$(b_0 + W) - b \le c$}{tag component with $\fR$; add face port}
}
$\iocount \gets |\{C : \fI \in \text{faces}(C) \text{ and } \fO \in \text{faces}(C)\}|$\;
\Return TileOperator\;
\end{algorithm}

\subsection{ISE ensemble evaluation}

\begin{algorithm}[H]
\caption{ISE$(k^2, R_{\min}, R_{\max}, W, H, M)$}\label{alg:ise}
\KwIn{Step bound $k^2$, radial range $[R_{\min}, R_{\max}]$, tile dimensions $W \times H$, strip count $M$}
\KwOut{$\fr$ profile for each shell}
$c \gets \lceil\sqrt{k^2}\rceil$\;
$\text{stride} \gets W + 2c$\;
$\text{offsets} \gets [c + j \cdot \text{stride} \text{ for } j = 0, \ldots, M{-}1]$\;
$\text{shells} \gets [(a_{\text{lo}}, a_{\text{hi}}) \text{ for } a_{\text{lo}} \in \{R_{\min}, R_{\min}{+}H, \ldots, R_{\max}{-}H\}]$\;
\ForEach{shell $(a_{\text{lo}}, a_{\text{hi}})$}{
  $\text{io\_counts} \gets [0] \times M$\;
  \For(\tcp*[f]{in parallel}){$j = 0$ \KwTo $M-1$}{
    $T \gets \text{ComputeTile}(a_{\text{lo}},\, \text{offsets}[j],\, W,\, a_{\text{hi}} - a_{\text{lo}},\, k^2)$\;
    $\text{io\_counts}[j] \gets \iocount(T)$\;
  }
  $\fr \gets |\{j : \text{io\_counts}[j] > 0\}| / M$\;
  \lIf{$\fr = 0$}{flag shell as moat candidate}
}
\Return $\fr$ profile\;
\end{algorithm}

\subsection{Compensated ISE}

\begin{algorithm}[H]
\caption{CompensatedISE$(k^2, R_{\text{target}}, W, H, M, N_{\text{shells}})$}\label{alg:compensated-ise}
\KwIn{Step bound $k^2$, target radius $R_{\text{target}}$, tile dimensions $W \times H$, strip count $M$, shells per probe $N_{\text{shells}}$}
\KwOut{$\fr$ profile}
$c \gets \lceil\sqrt{k^2}\rceil$\;
$\text{stride} \gets W + 2c$\;
$\text{offsets} \gets [c + j \cdot \text{stride} \text{ for } j = 0, \ldots, M{-}1]$\;
\For{$j = 0$ \KwTo $M-1$}{
  $b_{\text{center}} \gets \text{offsets}[j] + W/2$\;
  $a_{\text{base}}[j] \gets \lfloor\sqrt{R_{\text{target}}^2 - b_{\text{center}}^2}\rfloor$\;
}
\For{$n = 0$ \KwTo $N_{\text{shells}} - 1$}{
  $\text{io\_counts} \gets [0] \times M$\;
  \For(\tcp*[f]{in parallel}){$j = 0$ \KwTo $M-1$}{
    $a_{\text{lo}} \gets a_{\text{base}}[j] + n \cdot H$\;
    $T \gets \text{ComputeTile}(a_{\text{lo}},\, \text{offsets}[j],\, W,\, H,\, k^2)$\;
    $\text{io\_counts}[j] \gets \iocount(T)$\;
  }
  $\fr \gets |\{j : \text{io\_counts}[j] > 0\}| / M$\;
}
\Return $\fr$ profile\;
\end{algorithm}

\subsection{Fat-stripe campaign}

\begin{algorithm}[H]
\caption{FatStripe$(k^2, r_{\min}, r_{\max}, W, H, C)$}\label{alg:fat-stripe}
\KwIn{Step bound $k^2$, annular radii $[r_{\min}, r_{\max}]$, tile dimensions $W \times H$, column-chunk size~$C$}
\KwOut{$\texttt{blocked}$ (boolean), $\texttt{spanning\_count}$}
$c \gets \lceil\sqrt{k^2}\rceil$\;
$\text{sieve\_table} \gets \text{Precompute}(L = 110{,}000)$\;
$\text{composed} \gets \texttt{null}$\;
\For{$a_{\text{lo}} = r_{\min}$ \KwTo $r_{\max}$ \textbf{step} $H$}{
  $a_{\text{hi}} \gets \min(a_{\text{lo}} + H,\; r_{\max})$\;
  $b_{\max} \gets \min\!\big(\lceil\sqrt{r_{\max}^2 - a_{\text{lo}}^2}\rceil,\; a_{\text{hi}}\big)$\;
  $\text{stripe\_op} \gets \texttt{null}$\;
  \For{$b_{\text{lo}} = 0$ \KwTo $b_{\max}$ \textbf{step} $C \cdot W$}{
    $b_{\text{hi}} \gets \min(b_{\text{lo}} + C \cdot W,\; b_{\max})$\;
    \tcp{Phase 1: Sieve expanded chunk region (parallel over rows)}
    $\text{primes} \gets \text{SieveChunkRows}(a_{\text{lo}} - c,\; a_{\text{hi}} + c,\; b_{\text{lo}} - c,\; b_{\text{hi}} + c,\; \text{sieve\_table})$\;
    \tcp{Phase 2: Partition primes into virtual tiles with collar overlap}
    $\text{tile\_primes} \gets \text{Partition}(\text{primes},\; b_{\text{lo}},\; W,\; c)$\;
    \tcp{Phase 3: Sparse UF per tile $\to$ TileOperator (parallel)}
    \For(\tcp*[f]{in parallel}){each virtual tile $t$}{
      $\text{tile\_ops}[t] \gets \text{BuildTile}(a_{\text{lo}}, a_{\text{hi}}, b_t, b_t', k^2, \text{tile\_primes}[t])$\;
    }
    \tcp{Phase 4: Horizontal composition L$\to$R within chunk}
    $\text{chunk\_op} \gets \text{tile\_ops}[0]$\;
    \For{$t = 1$ \KwTo $|\text{tile\_ops}| - 1$}{
      $\text{chunk\_op} \gets \text{ComposeHorizontal}(\text{chunk\_op},\; \text{tile\_ops}[t],\; k^2)$\;
    }
    $\text{stripe\_op} \gets \text{ComposeHorizontal}(\text{stripe\_op},\; \text{chunk\_op},\; k^2)$\;
  }
  $\text{composed} \gets \text{ComposeVertical}(\text{composed},\; \text{stripe\_op},\; k^2)$\;
}
\tcp{Verdict: radius-based spanning check}
$\text{spanning\_count} \gets 0$\;
\ForEach{component~$C$ in \textup{composed}}{
  $\text{has\_inner} \gets \exists\, p \in C : \sqrt{p.a^2 + p.b^2} \le r_{\min} + c$\;
  $\text{has\_outer} \gets \exists\, p \in C : \sqrt{p.a^2 + p.b^2} \ge r_{\max} - c$\;
  \lIf{$\text{has\_inner} \wedge \text{has\_outer}$}{$\text{spanning\_count} \mathrel{+}= 1$}
}
$\texttt{blocked} \gets (\text{spanning\_count} = 0)$\;
\Return $\texttt{blocked}$, $\text{spanning\_count}$\;
\end{algorithm}


% ══════════════════════════════════════════════════════════════
\section{Feasibility Calculation Details}\label{app:feasibility}

\subsection{Neighbor counts by \texorpdfstring{$k^2$}{k\texttwosuperior}}

\begin{table}[ht]
\centering
\begin{tabular}{@{}cccc@{}}
\toprule
$k^2$ & Nonzero neighbor vectors & Collar $c$ & Max-excursion vector \\
\midrule
2  & 4   & 2 & $(1, 1)$ \\
26 & 80  & 6 & $(5, 1)$ \\
32 & 92  & 6 & $(4, 4)$ \\
36 & 112 & 6 & $(6, 0)$ \\
40 & 128 & 7 & $(6, 2)$ \\
\bottomrule
\end{tabular}
\caption{Neighbor counts and collar widths by step bound.}
\label{tab:neighbor-counts}
\end{table}

\subsection{Percolation threshold estimation for \texorpdfstring{$k^2 = 40$}{k\texttwosuperior\ = 40}}

Empirical degree-statistics measurements (via \texttt{fat-stripe --degree-stats}) confirm a universal critical threshold of $d_c \approx 4.0$ mean total degree across $k^2 = 32$, $36$, and~$40$.  At $k^2 = 36$, the percolation boundary ($\fr = 0$) was found at $R \approx 80.4\text{M}$ with measured mean total degree~$3.97$ (and~$3.92$ at the supercritical reference radius $R = 40\text{M}$ where $\fr > 0$).

Using $d_c = 4.0$ and total neighbor count $z = 128$ for $k^2 = 40$:
\[
  \rho_c(k^2 = 40) = \frac{d_c}{z} = \frac{4.0}{128} = 0.03125 = 3.125\%.
\]
Setting $\rho(R) = C / (2\ln R) = \rho_c$:
\[
  \ln R = \frac{C}{2\rho_c} = \frac{1.274}{2 \times 0.03125} = 20.38, \qquad R = e^{20.38} \approx 7.1 \times 10^8.
\]
This estimate of $R \approx 710\text{M}$ is broadly consistent with the empirical sigmoid midpoint $\Rhalf \approx 839\text{M}$ and the measured degree data, which shows the $d_c \approx 4.0$ crossing occurring in the range $R \approx 1.0\text{B}$--$1.2\text{B}$ (measured mean total degree: $4.03$ at $R = 1.05\text{B}$, $4.00$ at $R = 1.2\text{B}$).  The discrepancy between the $\rho$-based estimate and the measured crossing radius reflects the approximation in the prime density model $\rho(R) = C / (2\ln R)$.

\subsection{Memory requirements}

Per-tile working memory for $2000 \times 2000$ tile with $c = 7$:

\begin{table}[ht]
\centering
\begin{tabular}{@{}lr@{}}
\toprule
Component & Size \\
\midrule
Expanded tile bitmap $(2015^2)$            & $507\text{ KB}$ \\
Union-find parent array ($n$ \texttt{u32}) & ${\sim}\,16\text{ MB}$ \\
Union-find rank array ($n$ \texttt{u8})    & ${\sim}\,4\text{ MB}$ \\
Spatial hash                               & ${\sim}\,2\text{ MB}$ \\
\midrule
\textbf{Total per tile}                    & $\mathbf{{\sim}\,23\textbf{ MB}}$ \\
\bottomrule
\end{tabular}
\caption{Per-tile memory breakdown.}
\label{tab:memory}
\end{table}

With 6 parallel tiles (Jetson Orin Nano, 6 cores): ${\sim}\,138\text{ MB}$ peak.  With 64 parallel tiles (cloud): ${\sim}\,1.5\text{ GB}$.

\subsection{Tsuchimura comparison}

A Tsuchimura-style full sweep to $R = 1.5\text{B}$ at $k^2 = 40$ would require storing all primes in the disk of radius $1.5\text{B}$:
\[
  N_{\text{primes}} \approx \frac{\pi R^2 \cdot \rho(R)}{1} = \frac{\pi \cdot (1.5 \times 10^9)^2 \cdot 0.031}{1} \approx 2.2 \times 10^{17}.
\]
At 16 bytes per prime (coordinates + component ID): ${\sim}\,3.5\text{ PB}$.  Union-find adds comparable storage.  This is infeasible on any existing hardware.

The ISE, by contrast, processes one $2000 \times 2000$ tile at a time with ${\sim}\,23\text{ MB}$ working memory, achieving a memory reduction of approximately $10^{11}\times$.

\subsection{Primality validity ceiling}

The implementation uses a 9-witness deterministic Miller--Rabin test, valid for all $n < 3.825 \times 10^{18}$.  At radius~$R$, the maximum norm tested is approximately~$R^2$.  The hard ceiling is $R \le 1.955 \times 10^9$ ($\approx 1.955\text{B}$).  Beyond this, either additional witnesses or a BPSW test is required.  The $k^2 = 40$ campaign data through $R = 3.5\text{B}$ uses 12 witnesses for probes beyond the 9-witness ceiling.

\subsection{GPU throughput calculations}\label{app:gpu-throughput}

The throughput projections in Section~\ref{sec:ub-feasibility} are derived as follows.

\paragraph{Sequential per-tile time.}
On Jetson Orin Nano, the GPU primality kernel takes ${\sim}\,170\text{ ms}$ for a $2000 \times 2000$ tile.  Eliminating the CPU UF stage and D2H copy removes ${\sim}\,42\text{ ms}$, but the GPU UF kernel adds ${\sim}\,50\text{ ms}$ (estimate based on ECL-CC benchmarks on similar grid sizes), yielding ${\sim}\,180\text{ ms}$ total per tile.

\paragraph{Cross-GPU scaling.}
Scaling from Jetson to discrete GPUs uses the SM count and clock ratio:
\[
  t_{\text{GPU}} = t_{\text{Jetson}} \times \frac{\text{SM}_{\text{Jetson}}}{\text{SM}_{\text{GPU}}} \times \frac{f_{\text{Jetson}}}{f_{\text{GPU}}} \times \alpha,
\]
where $\alpha \approx 0.7$ accounts for memory bandwidth differences (Jetson uses shared LPDDR5; discrete GPUs use dedicated HBM/GDDR6X with higher bandwidth).

\paragraph{Batched throughput.}
With tile-per-SM batching, the effective per-tile time is the sequential time divided by the SM count, assuming sufficient tiles to fill all SMs:
\[
  t_{\text{batched}} = \frac{t_{\text{sequential}}}{\text{SM count}}.
\]
At $R = 1.5\text{B}$, the shell contains ${\sim}\,589\text{K}$ tiles---far exceeding any GPU's SM count---so the batching assumption holds.

\paragraph{Cost model.}
Cloud cost estimates use: RTX 4090 at ${\sim}\,\$2/\text{hr}$ (consumer cloud), A100 SXM4 at ${\sim}\,\$2/\text{hr}$ (reserved spot pricing).  Multi-GPU configurations assume linear scaling (tile-level parallelism with no inter-GPU communication).

% ══════════════════════════════════════════════════════════════
\begin{thebibliography}{9}

\bibitem{jaiganesh2018}
S.~Jaiganesh and M.~Burtscher,
\emph{A high-performance connected components implementation for GPUs},
in Proc.\ PPoPP, 2018, pp.~92--104.

\bibitem{sutton2018}
M.~Sutton, T.~Ben-Nun, and A.~Barak,
\emph{Optimizing parallel graph connectivity computation via subgraph sampling},
in Proc.\ SC, 2018, Article~14.

\end{thebibliography}

\end{document}
